\section{Background}
\label{sec:background}

The challenge of governing hierarchies of autonomous agents is not new, but it has become acute with the emergence of Large Language Model (LLM)-powered systems that can spawn sub-agents, delegate tools, and execute open-ended plans without predefined execution boundaries. This section surveys four intersecting research threads that collectively define the landscape into which the BX3 Framework's Recursive Spawning pillar enters: agent lifecycle management, accountability and provenance in distributed systems, fixed versus mutable delegation chains, and hierarchical task decomposition. The section closes by situating BX3 as a concrete architectural substrate that resolves tensions across these threads.

% ============================================================
\subsection{Agent Lifecycle Management in Multi-Agent Systems}
% ============================================================

A multi-agent system (MAS) is an ensemble of autonomous entities that coordinate to achieve goals beyond the capacity of any single agent. When those agents are LLM-powered, the lifecycle complexity multiplies: agents must be created, provisioned with capabilities, assigned tasks, monitored for performance, adapted when they fail, and retired cleanly when their function is complete. Recent literature has formalized this concern as a first-class design axis.

Moore's five-axis taxonomy for Hierarchical Multi-Agent Systems (HMAS) identifies lifecycle alignment as a core design dimension: higher layers govern long-horizon planning and accountability, while lower layers manage execution and local adaptation \cite{moore2025hmas}. This decomposition maps naturally to the spawning problem: when a parent agent subdivides a goal, each child agent enters a bounded lifecycle tied to its delegated task. The taxonomy further specifies that lifecycle stages (creation, assignment, execution, monitoring, adaptation, retirement) must be matched to appropriate hierarchy levels to preserve both autonomy and coherence \cite{moore2025hmas}.

The AgentOrchestra framework demonstrates this in practice: a conductor agent decomposes complex objectives into subtasks, assigns them to specialized sub-agents, and manages their lifecycle dynamically \cite{agentorchestra2025}. Critically, AgentOrchestra supports adaptive role reallocation — tasks can be repartitioned as conditions change — which requires lifecycle management at runtime rather than design time. This dynamic lifecycle property is essential for Recursive Spawning: a parent must be able to spawn a child, monitor its execution, and retire it without leaving orphaned processes or unresolved accountability gaps.

Agent lifecycle management also requires explicit resource governance. The Agent Contracts framework introduces resource-bounded contracts (tokens, API calls, temporal deadlines) as first-class primitives that govern lifecycle boundaries: when to spawn, when to transfer control, and when to terminate a child agent \cite{agentcontracts2025}. This contract-based lifecycle enforcement provides a structural model for the BX3 Worksheet mechanism, which similarly deploys containerized logic sets with defined execution boundaries.

The HASHIRU system extends this further with an autonomous cost model that governs hiring and firing of agents based on budget and performance constraints \cite{hashiru2025}. This represents a concrete implementation of lifecycle-as-resource-management: agents are not merely spawned but instantiated on demand, evaluated continuously, and dissolved when their marginal value falls below cost. BX3's Recursive Spawning adopts this principle at the architectural level — each spawned child carries a hard-coded parent pointer and bounded execution scope, making lifecycle management a structural guarantee rather than a policy choice.

% ============================================================
\subsection{Accountability and Provenance in Distributed Systems}
% ============================================================

When a chain of agents collectively produces an outcome, the question of who is accountable for each decision point becomes non-trivial. Accountability in distributed AI systems has emerged as a central concern for both regulatory bodies and technical architects. The EU AI Act and NIST's AI Risk Management Framework now demand defined roles, decision provenance, and human oversight for high-risk automated decisions \cite{dzone2025accountability}. Meeting these requirements architecturally — rather than as post-hoc compliance overlays — demands provenance systems that capture every delegation, every tool call, and every human authorization in an immutable, auditable chain.

The Human Delegation Provenance (HDP) Protocol represents the most rigorous approach to this problem: a lightweight, cryptographic protocol that binds human authorization to a session, and then records every subsequent agent delegation as a signed, append-only hop in a provenance chain \cite{hdp2026, hdp2025ietf}. Each hop is cryptographically signed using Ed25519, and the entire chain can be verified offline using only the issuer's public key and the session identifier — no registry lookups or trusted third parties required. HDP addresses a critical gap in prior standards (OAuth 2.0 Token Exchange, JWT, UCAN), which lack multi-hop, append-only human-provenance capabilities \cite{hdp2025arxiv}.

PROV-AGENT extends the W3C PROV standard with agent-centric metadata — prompts, responses, and decision rationale — linked via the Model Context Protocol (MCP) to end-to-end workflow provenance \cite{provcagent2025}. This enables near real-time provenance tracking across federated, heterogeneous agent deployments. The framework was evaluated across edge, cloud, and HPC environments, demonstrating that fine-grained agent provenance is queryable at scale without degrading system performance. For BX3, PROV-AGENT's provenance model provides a structural template: each BX3 loop emits provenance records that link the Purpose Layer's intent through the Bounds Engine's reasoning to the Fact Layer's enforcement, creating a complete audit trail from authorization to execution.

TrustTrack proposes a complementary four-layer protocol stack (Cloud → AI → Agent → Blockchain) that embeds cryptographic guarantees directly into agent infrastructure \cite{trusttrack2025}. Key properties include DID-based agent identity, signed actions, auditable traces, decentralized policy registration, and tamper-resistant behavioral logs. TrustTrack frames compliance as a design constraint rather than an external audit, achieving provenance through architectural embedding rather than retrofitted logging.

AIAuditTrack introduces dynamic interaction graphs where AI entities are nodes and time-specific edges capture behavioral trajectories, enabling cross-system supervision and traceability \cite{aaudittrack2025}. A risk-diffusion algorithm traces the origin of risky behaviors and propagates warnings across affected entities. This graph-based accountability model is orthogonal to but compatible with BX3's chain structure: each BX3 loop can emit events into a graph while maintaining its own immutable parent-pointer chain.

The Chain of Consciousness (CoC) protocol proposes a two-layer cryptographic approach to agent provenance, using an append-only SHA-256 hash chain of lifecycle events anchored to Bitcoin via OpenTimestamps \cite{chainofconsciousness2026}. Layer 2 extends provenance to fleets, including inter-agent delegation records. CoC introduces a governance mechanism — Proof of Continuity — where votes and influence derive from verified chain length rather than capital or computation. This positions agent age and continuous operation as trust primitives, a concept that BX3's immutable parent chains reinforce structurally.

ProML introduces a decentralized provenance management platform for distributed ML systems using blockchain and smart contracts, creating a tamper-evident trail of how ML assets are created, modified, and transferred across trust boundaries \cite{proml2022}. The Artefact-as-a-State-Machine model maps ML artifacts to blockchain transactions, enabling immutable audit records. This blockchain-native approach to provenance provides a deployment pattern for BX3 systems operating in high-assurance environments where regulatory auditability is non-negotiable.

% ============================================================
\subsection{Fixed Versus Mutable Delegation Chains}
% ============================================================

Delegation chains in multi-agent systems can be classified along a spectrum from fixed to mutable. Fixed delegation chains pre-define the complete agent hierarchy at design time: parent-child relationships are static, and no agent can spawn new agents beyond those explicitly specified. Mutable delegation chains permit runtime modification of the agent topology — agents can spawn new agents, re-delegate tasks, and restructure the hierarchy dynamically.

Fixed delegation provides predictability and auditability: the complete chain is known in advance, enabling deterministic accountability mapping. However, fixed delegation fails under conditions of task complexity, environmental variability, or resource scarcity — precisely the conditions that motivate recursive spawning. A static hierarchy cannot adapt when a spawned sub-task exceeds its parent's capability scope or when a child agent encounters an unanticipated failure mode requiring escalation.

Mutable delegation enables adaptive problem-solving but introduces accountability fragmentation. When an agent spawns a child mid-execution, the provenance chain must be extended dynamically — and if the parent subsequently crashes, is compromised, or exceeds its authority, the delegation chain may become unverifiable or orphaned. The Delegation Chain Calculus (DCC) formalized by SentinelAgent defines seven properties for secure, auditable delegation, including authority narrowing, policy preservation, forensic reconstructibility, cascade containment, scope-action conformance, output schema conformance, and intent preservation \cite{sentinelagent2026}. These properties constitute a formal specification for what mutable delegation chains must satisfy to remain accountable: a mutable chain that violates any of these properties is, by definition, no longer trustworthy.

BX3's Recursive Spawning resolves this tension architecturally. Each spawned child receives an immutable pointer to its parent's Purpose — not to the parent's Bounds Engine or Fact Layer, but specifically to the accountability-anchoring Purpose function. This pointer survives any mutation of the agent topology: even if the parent is dissolved, retired, or compromised, the child's accountability chain remains anchored to the parent's original intent. The chain is mutable in topology (agents can spawn, recombine, and retire) but immutable in its accountability anchoring (every child always traces back to a named, accountable Purpose). This is the key distinction that separates BX3's approach from both static delegation (which is safe but inflexible) and naive mutable delegation (which is flexible but unaccountable).

% ============================================================
\subsection{Hierarchical Task Decomposition in AI}
% ============================================================

Hierarchical task decomposition — the process by which a complex goal is recursively broken into simpler sub-goals handled by specialized agents — is a foundational technique in AI planning and multi-agent orchestration. The fundamental insight is that no single agent, human or artificial, can optimally decompose and execute a complex, multi-horizon plan in a single reasoning pass. Decomposition enables focused specialization, parallel execution, and bounded complexity management.

ReAcTree introduces a tree-structured approach to long-horizon task planning where a root agent decomposes goals into subgoals that spawn new agent nodes dynamically \cite{reactree2025}. Each agent node can reason, act, and further expand the tree, creating a recursive, parent-child-like relationship as the task unfolds. Control-flow nodes coordinate execution strategies across the tree, managing how subgoals are assigned, pursued, and when to spawn further nodes. Critically, ReAcTree uses episodic memory (goal-specific examples) and working memory (environment-specific observations) to support hierarchical decision-making — ensuring that spawned agents retain contextual coherence with their parent's intent.

AgentSpawn extends this model with the Spawn-Resume Protocol: seamless transfer of agent state and memory when spawning child agents, enabling continuity across task boundaries \cite{agentspawn2025}. An adaptive spawning policy uses runtime complexity metrics to decide when to spawn or reallocate agents, addressing unforeseen task difficulty dynamically. A Memory Coherence Manager ensures consistent, conflict-free concurrent edits when multiple child agents modify overlapping regions. AgentSpawn demonstrates ~34\% higher task completion versus static baselines and ~42\% reduction in memory overhead through selective context slicing, validating that hierarchical decomposition with memory transfer outperforms both flat and rigidly hierarchical approaches.

FIRMHIVE applies hierarchical decomposition to firmware security analysis through a runtime Tree of Agents (ToA) \cite{firmhive2025}. Tasks are recursively broken down into executable primitives delegated to specialized agents, with delegation treated as an actionable primitive rather than a static handoff. The ToA enables decentralized coordination where agents collaboratively reason across multiple files, components, and analysis stages — demonstrating that hierarchical decomposition scales to complex, heterogeneous domains without centralized bottlenecks.

The Recursive Delegation Engine (RDE) formalizes the decision process at each decomposition node: a parent agent decides whether to execute, decompose further, or delegate to a child agent based on runtime complexity signals \cite{rde2025}. The delegation topology is modeled as a directed graph G = (A, E), where nodes are agents and edges indicate possible delegation paths. Decision methods include epsilon-greedy exploration, UCB/Beta-UCB multi-armed bandit approaches, and Thompson Sampling — enabling learning-based trust optimization across the delegation chain. The utility function guiding decomposition decisions combines execution cost and contextual constraints: U(instr) = $\alpha(1 - c) + \beta (L / L_{max})$, where $c$ is cost and $L$ is current context window utilization \cite{rde2025}.

BX3's Recursive Spawning adopts the structural intuition from this body of work — that hierarchical decomposition is necessary for scalable problem-solving — while imposing a stronger accountability constraint: every decomposition event must produce a child loop with an immutable parent Purpose pointer. This prevents the common failure mode in hierarchical task decomposition where a deeply spawned agent loses coherence with the original intent, or where accountability becomes diffuse across multiple levels of delegation.

% ============================================================
\subsection{The BX3 Framework as Substrate for Immutable Parent Chains}
% ============================================================

The research surveyed above reveals a field that has independently identified the structural needs for hierarchical, accountable, lifecycle-managed agent systems — but has not yet synthesized these needs into a single coherent architectural model that enforces accountability as a non-negotiable property at every node.

HMAS taxonomy provides the structural language for lifecycle alignment and hierarchy design \cite{moore2025hmas}. AgentOrchestra, AgentSpawn, and FIRMHIVE provide design patterns for dynamic decomposition and spawning \cite{agentorchestra2025, agentspawn2025, firmhive2025}. HDP, PROV-AGENT, TrustTrack, and AIAuditTrack provide provenance and accountability mechanisms \cite{hdp2026, hdp2025ietf, provcagent2025, trusttrack2025, aaudittrack2025}. Agent Contracts and the RDE provide resource-aware lifecycle management \cite{agentcontracts2025, rde2025}. The Chain of Consciousness adds cryptographic anchoring and proof-of-continuity governance \cite{chainofconsciousness2026}.

What none of these frameworks individually provides is an architectural guarantee that every spawned child, regardless of how deep the hierarchy grows, maintains an immutable, verifiable chain back to a named Purpose — and that unresolved escalations always fail upward into human accountability, never downward into algorithmic drift. This is the gap the BX3 Framework's Recursive Spawning pillar is designed to fill.

The BX3 approach treats the parent pointer not as a soft reference but as a structural guarantee: the Worksheet deployed to each child contains a cryptographically signed, immutable pointer to the parent node's Purpose function. This pointer is embedded at the architectural level, not added as metadata. The result is a system where delegation chains are simultaneously mutable in topology (agents can spawn, adapt, and restructure) and immutable in accountability (every action traces to a named Purpose, and every unresolved state escalates to the human who set that Purpose). This dual property — mutable topology, immutable accountability — is what distinguishes BX3's Recursive Spawning from both fixed delegation and from mutable delegation systems that lack accountable anchoring.

\end{document}